Large language models can produce answers that are fluent and plausible while
still being unsupported by the available evidence.
For this project, the main concern is not hallucination as a vague catch-all
label, but unsupported answering delivered as if it were reliable.
That framing matters because improving average accuracy does not automatically
reduce overconfident unsupported outputs.

This paper takes a deliberately narrow starting point.
Instead of attempting to solve hallucination in general, we study a controlled
evidence-grounded question answering task in which a system must either answer
or abstain.
The current implementation uses SQuAD~2.0 \citep{rajpurkar2018know}
with one question and one evidence passage as input.
Stage 1 compares two training regimes:
\begin{enumerate}[leftmargin=1.25em]
  \item a baseline trained only on answerable examples, and
  \item an abstain-aware variant trained to treat unsupported cases as a
  first-class target behavior.
\end{enumerate}
Stage 2 then adds a verifier over question, passage, and predicted answer in
order to estimate whether an answer is actually supported before it is shown to
the user.

The central hypothesis is simple:
explicit abstention supervision should reduce unsupported answers without
causing an unacceptable collapse in answer quality on answerable questions.
That hypothesis is narrow enough to evaluate cleanly, but still meaningful as a
first step toward a broader program on grounded answering, support verification,
and calibrated uncertainty.
The next immediate question is whether a learned support signal can recover some
of the remaining quality-versus-grounding gap without undoing the abstention
gain from Stage 1.

The intended application setting is also narrow.
\system\ is most naturally suited to bounded, evidence-grounded QA scenarios in
which answers should be tied to supplied documents and abstention is preferable
to unsupported guessing, such as internal knowledge assistants, support systems
over approved documentation, and policy or compliance QA.
It is not intended as a general solution for open-domain truthfulness, balanced
opinion generation, or fairness beyond the scope of the provided evidence.

This framing is also consistent with recent work arguing that hallucination
persists partly because modern evaluation pipelines reward guessing more than
appropriate uncertainty \citep{kalai2025why}.
Our contribution is narrower and more operational:
rather than explaining hallucination in general, Stage 1 tests whether an
evidence-grounded QA system can be trained and evaluated so that abstention is
treated as a useful target behavior instead of a failure by default, and
Stage 2 tests whether a learned support verifier can improve downstream
answer-versus-abstain decisions on top of that base system.

This LaTeX scaffold is intentionally written as a living project paper.
Some sections already reflect stable design decisions, while others are
structured to absorb final metrics, figures, and ablations once the
experimental results are locked in.
